% 原文 p.81
\hypertarget{srcpage:81}{}\label{srcpage:81}
\subsubsectionlink{デコヒーレンスの量子極限における指針状態}
次に、系の発展と比べて環境のモードが「遅い」状況を考えよう。すなわち、環境が利用できる最高振動数、つまりエネルギーが、系のエネルギー固有状態間の間隔より小さいと仮定する。この場合、環境が監視できるのは運動の定数である量だけである。縮退のない場合、その量は系のエネルギーなので、環境誘起超選択によって系のエネルギー固有状態、すなわち系の自己Hamiltonianの固有状態が選ばれる。この過程は、PazとZurek\cite{ref103}が初めて詳しく調べ、デコヒーレンスの量子極限と呼んだ。

量子力学の入門教科書では、閉じた系の、そのようなエネルギー固有状態には、Hamiltonianの作用のもとで定常であるため、特別な役割を与えるのが普通である。実際、その後に扱う問題の多くは、水素原子などのさまざまな微視的物理系を記述する、異なるHamiltonianのエネルギー固有状態の計算に関わる。しかし重ね合わせ原理によれば、それでもなお、これらのエネルギー固有状態の任意の重ね合わせは、完全に正当であるはずである。

したがって、系の自己Hamiltonianが発展を支配する極限でのエネルギー固有状態の選択は、環境の存在をまったく無視してよい状況、すなわち教科書でよく扱う上述の閉じた系の場合とは、同じでないことを強調するのが、非常に重要である。むしろ環境は、系のエネルギーを連続的に監視し、系のエネルギー固有基底でデコヒーレンスを起こし、エネルギー固有状態の重ね合わせを局所的に抑制する、決定的な役割を果たす。したがって指針状態の形は、系の自己Hamiltonianの固有状態に支配されるが、観測可能な重ね合わせの空間に実効的超選択則を課すのは、ここでも環境の存在なのである。

\subsubsectionlink{指針状態を決める一般的方法}
関心のある現実的な場合の大半では、可換性基準(2.89)は、普通、近似的にしか満たせない\cite{ref084,ref086}。たとえば2.5.4節で述べたように、位置演算子の個々の固有状態は、物理的物体の正当な量子状態ではない。相互作用Hamiltonian(2.91)と可換性基準(2.89)によれば、その固有状態を系の優先指針状態に選ぶはずだが、実際の優先状態は、狭い位置空間波束となる。つまり、位置空間に有限の広がりを持つ、位置の近似的固有状態だけを表す。他の状況では、可換性基準(2.89)を厳密に満たす観測可能量を、そもそも見いだせないかもしれない。

% 原文 p.82
\hypertarget{srcpage:82}{}\label{srcpage:82}
したがって、相互作用Hamiltonianの対角化や可換性基準(2.89)の適用は、量子測定極限の成り立つ単純な場合に、優先状態と観測可能量を決める、概念的に直観的な方法であるが、より複雑な状況では、適切で実行可能であるとは限らない。そのため、環境中に浸されても完全に影響されない厳密な指針状態を見いだすという目標は、しばしば環境との相互作用に対して最も頑健な系の状態を決めるという目標へ弱められる。すると近似的な指針状態が、環境との不可避の相互作用に最小限しか影響されない準古典状態の、最善の表現になる。また、関心のある現実的な物理状況の多くでは、系の自己Hamiltonianと、系と環境の相互作用Hamiltonianが、おおむね等しい強さで寄与する。たとえば5.2節の量子Brown運動模型ではそうである。その場合、内在的力学を無視する量子測定極限も、「遅い」環境に対応するデコヒーレンスの量子極限も、当てはまらない。

したがって、環境との相互作用が選ぶ優先状態を決めるには、より一般的な方法が必要である。そうした方法は、Zurek\cite{ref084,ref104}、およびZurek、Habib、Paz\cite{ref086}が、予測可能性のふるいという方策の名のもとに発展させた。基本的な考えは、縮約密度行列の純度やvon Neumannエントロピー（2.4.3節）など、系に導入されるデコヒーレンスの量に対する適切な尺度を用いることである。そして、系と環境の全Hamiltonianのもとで発展する、系の初期状態の大きな集合について、純度またはvon Neumannエントロピー、あるいは選んだ他の尺度の時間依存性を計算する。デコヒーレンスに最も影響されない状態は、純度の減少やvon Neumannエントロピーの増加が最小のものになる。この方法を適用すると、可能な優先状態を、その「古典性」、すなわち環境との相互作用に対する頑健性に従って、順位付けできる。

このように、Hilbert空間を「ふるいにかけ」、環境との相互作用に対する頑健性によって状態を分類すると考える。最も安定な状態は、最も予測しやすい状態でもある。これがZurekの「予測可能性のふるい」という用語の動機である。これらの状態で純度の喪失やvon Neumannエントロピーの増加が最小になることは、系の状態についての「情報」の喪失も最小となることを意味するからである。von Neumannエントロピーが零の純粋状態は、系の状態についての最大の知識に対応することを思い出そう。

一般に、予測可能性のふるいが選ぶ優先状態は、状態の頑健性の順位付けに、どの尺度を使うかによって異なりうる。すでに純度とvon Neumannエントロピーの二つを挙げたが、デコヒーレンスと頑健性については、他の意味のある尺度も存在し、今後さらに提案される可能性もある。

% 原文 p.83
\hypertarget{srcpage:83}{}\label{srcpage:83}
特定の模型では、予測可能性のふるいが選ぶ状態は、尺度の具体的な選択にほとんど左右されないことが、明示的に示されている。たとえば量子Brown運動模型（5.2節）では、異なる尺度が同じ最小不確定性波束に至る\cite{ref006,ref017,ref084,ref105,ref106}。少なくとも巨視的極限では、異なる基準から得られる安定した指針状態は、とても似たものになると予想するのが妥当である（Zurekの議論\cite{ref016}も参照）。

\subsubsectionlink{準古典的性質の選択}
優先状態の環境誘起選択についての結果をまとめ、2.5.2節で示した一般的な優先基底の問題、および周囲の世界で特定の物理量だけを観測することへの含意を論じよう。

系と環境の相互作用Hamiltonianは、周囲の光子や空気分子などの粒子が、調べている系と相互作用する散乱過程を記述することが多い。このような過程を記述する力の法則は、Newtonの力やCoulomb力の$\propto r^{-2}$のように、通常、距離の何らかのべきに依存するため、相互作用Hamiltonianは、普通、位置演算子と可換である。したがって可換性条件(2.89)によれば、指針状態は位置の近似的固有状態となる。位置が典型的に経験上確定した性質となることは、このように、大半の相互作用が距離に依存することによって説明できる。量子から古典への移行で位置が果たす特別な役割の、この起源を、初めて明瞭に指摘して分析したのはZurek\cite{ref008,ref009,ref013}であった。その後、JoosとZehの散乱模型\cite{ref007}は、周囲の光子や空気分子が、塵粒子などの小さな物体の空間構造を、どのように連続的に測定し、狭い位置空間波束の非真正混合への急速なデコヒーレンスをもたらすかを、直接示した（第3章）。

通常はエネルギー固有状態にある微視的系でも（下記参照）、周囲の媒質に強く結合する、特徴的な空間構造を持つ場合には、同様の結果が成り立つことがある。たとえば糖のようなキラル分子は、異なるエネルギー固有状態の重ね合わせである、左手型と右手型のキラリティ固有状態に、常に観測される\cite{ref107,ref108}（文献\cite{ref017}の3.2.4節も参照）。ここでも、これらの分子の異なる空間構造を、環境が散乱過程を通じて連続的に監視することによって説明できる。この環境による監視は、たとえばパリティやエネルギーの測定を目的とする測定装置で、通常達成できるよりはるかに強い「外界」との結合を生む。さらに、これらの分子をエネルギー固有状態に準備しようとしても、環境に対して安定なキラリティ固有状態へ、直ちにデコヒーレンスを起こし、位置が優先される量に選ばれる。

% 原文 p.84
\hypertarget{srcpage:84}{}\label{srcpage:84}
他方、原子などの微視的系は、相互作用Hamiltonianがエネルギー以外の観測可能量に依存する場合でも、通常、エネルギー固有状態に見いだされることがよく知られている。これは、そのような場合、系の発展が自己Hamiltonianに支配され、環境の最高振動数が系のエネルギー固有状態間隔より小さいという意味で、環境が比較的遅いことに注意すれば、容易に説明できる。2.8.2節で論じたように、このデコヒーレンスの量子極限\cite{ref103}は、エネルギー固有状態の環境誘起選択をもたらす。系のエネルギーを環境が監視するため、異なるエネルギー固有状態間の干渉が、絶えず抑制される。

もう一つ興味深い例は、電荷の環境誘起超選択である。自然界では、電荷演算子の固有状態だけが観測され、異なる電荷の重ね合わせは決して観測されない。対応する超選択則は、初めは仮定されるだけだったが\cite{ref078,ref079}、後にGiulini、Kiefer、Zeh\cite{ref085}が、電荷と、その電荷自身のCoulomb遠方場との相互作用によって説明した（文献\cite{ref083}も参照）。この場は環境として働き、電荷の重ね合わせを、電荷固有状態の非真正混合へ直ちにデコヒーレンスさせる。

まとめると、環境との相互作用から現れる、優先指針状態の型について、異なる三つの場合を区別した。
\begin{enumerate}
\item 量子測定極限。系の発展が$\hat H_{\mathrm{int}}$、すなわち環境との相互作用に支配される場合、優先状態は$\hat H_{\mathrm{int}}$の固有状態、したがってしばしば位置固有状態となる。
\item デコヒーレンスの量子極限。環境が遅く、系の自己Hamiltonian$\hat H_{\mathcal S}$が発展を支配する場合、優先状態はエネルギー固有状態、すなわち$\hat H_{\mathcal S}$の固有状態となる\cite{ref103}。これは微視的領域でよく起こる。
\item 中間領域。系の発展が$\hat H_{\mathrm{int}}$と$\hat H_{\mathcal S}$に、おおむね等しい強さで支配される場合、優先状態は前の二つの場合の折衷となる。たとえば、よく研究される量子Brown運動模型（5.2節）では、相互作用Hamiltonian$\hat H_{\mathrm{int}}$は系の位置の監視を記述する。しかし、$\hat H_{\mathcal S}$が生む内在的な力学を通じ、この監視は運動量の間接的なデコヒーレンスももたらす。$\hat H_{\mathrm{int}}$と$\hat H_{\mathcal S}$の組合わさった影響によって、位相空間、すなわち位置と運動量の両方に局在した優先状態が現れる（5.2.5節、および文献\cite{ref016,ref017,ref086,ref106,ref109}）。
\end{enumerate}
実際、この三領域は、優先指針状態は、全Hamiltonianが生成する発展のもとで、最も頑健な系の状態であるという、一般的規則の極限としても理解できる。全Hamiltonianのどの項が、相対的なエネルギー尺度で支配的かに応じて、指針基底は、相互作用Hamiltonianの固有状態、自己Hamiltonianの固有状態、あるいは両者の折衷として生じる状態によって構成される。

% 原文 p.85
\hypertarget{srcpage:85}{}\label{srcpage:85}
優先基底の問題に対する環境誘起超選択の方法の明確な長所は、単に測定記録を確定させたり、通常どの物理量、たとえば位置が確定していると知覚されるかという経験に合わせたりする、その場しのぎの方法で優先基底を選ばない点にある。代わりに選択は、物理的で観測者に依存しない根拠、すなわち系と環境の相互作用Hamiltonianの構造によって動機付けられる。可能な量子力学的重ね合わせの広大な空間が、これほど小さくなるのは、物理的相互作用を支配する法則が、位置、運動量、電荷など、少数の物理量だけに依存するからである。そして、まさにそれらが確定した性質として私たちに現れることは、優先基底が相互作用の形に依存することによって説明される。したがって古典性の出現は、系と環境の相互作用を支配する物理法則の構造に根差している。
